\documentclass{article}

\usepackage[final]{neurips_2025}

\usepackage[utf8]{inputenc}
\usepackage[T1]{fontenc}
\usepackage{hyperref}
\usepackage{url}
\usepackage{booktabs}
\usepackage{amsfonts}
\usepackage{amsmath}
\usepackage{amssymb}
\usepackage{nicefrac}
\usepackage{microtype}
\usepackage{xcolor}
\usepackage{graphicx}
\usepackage{algorithm}
\usepackage{algorithmic}
\usepackage{multirow}
\usepackage{enumitem}
\usepackage{array}

\newcommand{\system}{\textsc{GovernedKG}}
\newcommand{\eg}{\textit{e.g.}}
\newcommand{\ie}{\textit{i.e.}}
\newcommand{\etal}{\textit{et al.}}

\title{Governed Knowledge Graphs:\\Domain-Owned Subgraphs for Knowledge Creation, Governance, and Application}

\author{
  Pranav Bykampadi \\
}

\begin{document}

\maketitle

\begin{abstract}
We present \system{}, a governed knowledge graph architecture in which extracted knowledge is not treated as an unowned collection of atomic triples, but as a set of domain-owned subgraphs managed by expert agents. The core object is a \emph{Governed Knowledge Graph}, formally defined as $G=(E,T,D,\phi,\gamma)$, where $E$ is a set of entities, $T \subseteq E \times R \times E$ is a set of triples, $D$ is a set of domains, $\phi$ maps entities to their governing domains, and $\gamma$ is the governance function that decides whether a proposed triple is approved, rejected, revised, or escalated before it enters memory. The creation pipeline processes documents through specialized extraction agents, uses blackboard deliberation for low-confidence hypotheses, bootstraps an organizational chart during creation, and routes every admitted triple through a governance protocol with audit logging. On SciERC, governed creation lifts triple F1 by $+271\%$ relative at 10 documents and $+23.5\%$ relative at 50 documents, while reducing triple hallucination by 7--25 pp, at a $4.55\%$ wall-clock overhead. At 100 documents the governance layer remains essentially perfect (routing exact match $1.000$, governance completeness $1.000$, audit-trail integrity $1.000$, domain coverage $0.984$) even as extraction noise grows. A strict review benchmark achieves $0\%$ false accepts with $100\%$ rejection of corrupted triples. On downstream QA on the same 50-document governed KG, a four-way comparison against a flat-RAG baseline, a GraphRAG-style community-summary baseline, and our governed router shows that all configurations achieve high faithfulness ($\geq 0.98$) once correct abstention on negative questions is properly credited; the full governed stack (\texttt{domain\_advanced}) is the only system at full coverage ($1.000$) with the lowest hallucination rate ($0.9\%$) among the governance configurations. These results support the claim that the primary contribution is not state-of-the-art extraction, but a new governed memory structure for knowledge creation, organization, and downstream use.
\end{abstract}

\section{Introduction}

Most knowledge graph pipelines optimize for extraction: given documents, produce entities and relations as accurately as possible. In this framing, knowledge is treated as a set of atomic facts, and the graph is primarily a storage target. Our work starts from a different question: \emph{how should machine knowledge be organized and admitted into memory when knowledge is owned, reviewed, and propagated by domain-specialized agents?}

The central hypothesis of this paper is that knowledge should be represented as a governed structure rather than a passive fact store. A fact should not merely be extracted and appended; it should be routed to the domain expert agents that own the relevant subgraph, reviewed under an explicit governance protocol, and admitted with an audit trail. In this view, information is not an isolated triple. It is a local subgraph containing neighboring entities, relation patterns, provenance, domain ownership, and conflict context.

This reframing shifts the contribution away from question answering as the main object. Question answering is only one application on top of the governed structure. Knowledge ingestion, update review, conflict handling, and organizational structure are equally first-class. We therefore center the paper on three contributions:

\begin{enumerate}[nosep, leftmargin=*]
    \item a formal data structure for governed machine knowledge,
    \item a creation pipeline in which governance exists during KG construction rather than only after the graph is built,
    \item empirical governance results showing that routing, coverage, auditability, and accept/reject behavior are measurable and strong.
\end{enumerate}

\section{Governed Knowledge Graph}
\label{sec:definition}

\subsection{Formal Definition}

\textbf{Definition 1.} A \emph{Governed Knowledge Graph} is a tuple
\begin{equation}
G = (E, T, D, \phi, \gamma)
\end{equation}
where:
\begin{itemize}[nosep, leftmargin=*]
    \item $E$ is a set of entities,
    \item $T \subseteq E \times R \times E$ is a set of typed triples,
    \item $D = \{d_1, \dots, d_k\}$ is a set of domains,
    \item $\phi : E \rightarrow 2^D$ is an ownership function assigning each entity to one or more governing domains,
    \item $\gamma : T \rightarrow \{\texttt{approve}, \texttt{reject}, \texttt{revise}, \texttt{escalate}, \texttt{auto\_approve}\}$ is a governance function over proposed triples.
\end{itemize}

This definition differs from a standard knowledge graph in two ways. First, the graph contains an explicit organizational layer through $D$ and $\phi$. Second, admission into the triple set is mediated through $\gamma$, not direct insertion. The implementation mirrors this definition directly through a two-phase protocol:
\begin{equation}
\texttt{propose\_triple} \;\rightarrow\; \texttt{GovernanceDecision} \;\rightarrow\; \texttt{commit\_decision}.
\end{equation}

\subsection{Why This Is a New Data Structure}

In an ordinary KG, a triple either exists or does not exist. In a governed KG, each triple also has:
\begin{itemize}[nosep, leftmargin=*]
    \item an ownership assignment,
    \item a governance action,
    \item a rationale,
    \item an audit record,
    \item potentially a revision or escalation path.
\end{itemize}

This changes the graph from a passive fact store into a managed memory structure. The main research idea is therefore not ``multi-agent extraction'' in isolation, but \emph{knowledge creation and proliferation under domain-governed memory}.

\section{System Overview}
\label{sec:overview}

The system has three layers:
\begin{enumerate}[nosep, leftmargin=*]
    \item \textbf{Creation layer}: turns documents into entities and triples.
    \item \textbf{Governance layer}: organizes knowledge into domains and controls admission into memory.
    \item \textbf{Application layer}: uses the governed graph for downstream tasks such as question answering and incremental enrichment.
\end{enumerate}

The contribution of the paper lives in the first two layers. The third layer demonstrates usefulness but is not the center of the claim.

\section{Knowledge Graph Creation Workflow}
\label{sec:creation}

\subsection{Input and Orchestration}

The system receives a corpus of documents. A top-level orchestrator agent manages the end-to-end creation process over the full corpus and over each document instance. The pipeline is mostly sequential: each specialized agent transforms or validates the current representation and passes it to the next stage. This is intentionally simple. The system only activates deliberation when an agent signals uncertainty or conflict.

\subsection{Pipeline Steps}

The document-to-graph workflow is:

\begin{algorithm}[t]
\caption{Governed KG Creation}
\label{alg:gkg-creation}
\begin{algorithmic}[1]
\REQUIRE Corpus $\mathcal{C}$, Governed KG $G=(E,T,D,\phi,\gamma)$
\FOR{each document $x \in \mathcal{C}$}
    \STATE $\text{segments} \leftarrow \textsc{DocumentProcessor}(x)$
    \STATE $\text{schema} \leftarrow \textsc{DomainClassifier}(\text{segments})$
    \STATE $D_{\text{boot}} \leftarrow \textsc{BootstrapDomains}(\text{schema})$
    \STATE $\text{entities} \leftarrow \textsc{EntityExtractor}(\text{segments}, \text{schema})$
    \STATE $\phi \leftarrow \textsc{AssignEntitiesToDomains}(\text{entities}, D_{\text{boot}})$
    \STATE $\text{triples} \leftarrow \textsc{RelationExtractor}(\text{segments}, \text{entities}, \text{schema})$
    \STATE $\text{linked} \leftarrow \textsc{EvidenceLinker}(\text{triples}, \text{segments})$
    \STATE $\text{verified} \leftarrow \textsc{VerificationAgent}(\text{entities}, \text{linked}, x)$
    \FOR{each triple $t$ in $\text{verified}$}
        \STATE $\delta \leftarrow \textsc{propose\_triple}(t)$
        \STATE $\textsc{commit\_decision}(\delta)$
    \ENDFOR
\ENDFOR
\RETURN Governed KG $G$
\end{algorithmic}
\end{algorithm}

We now describe each stage in the same order that knowledge moves through the system.

\paragraph{1. DocumentProcessor.} The document is segmented into coherent chunks. This stage normalizes text, preserves metadata, and produces the segment sequence consumed by downstream agents.

\paragraph{2. DomainClassifier.} The classifier derives a working schema for the document. In benchmark mode it uses a fixed schema; in open-world mode it discovers entity and relation types from the content. This schema is the system's first coarse view of what kind of knowledge the document contains.

\paragraph{3. BootstrapDomains.} Governance begins during creation, not after the fact. From the discovered schema, the system constructs preliminary domains before triple admission. These domains are lightweight ownership shells that later grow into full subgraphs.

\paragraph{4. EntityExtractor.} The extractor identifies entity mentions, canonical labels, and candidate types. It is instructed to favor the shortest meaningful noun phrase and to preserve short scientific terms, abbreviations, tool names, and task phrases rather than discarding them as generic text.

\paragraph{5. AssignEntitiesToDomains.} Once entities exist, the ownership function $\phi$ becomes partially instantiated. Each entity is assigned to one or more domains using type, token, and topic compatibility. This is what allows the system to say, during creation, who owns a proposed fact.

\paragraph{6. RelationExtractor.} The extractor identifies relations between the current entities and produces candidate triples. In fixed-schema benchmark mode, the extractor is constrained to the benchmark label set; in open-world mode, it can propose new relation structure.

\paragraph{7. EvidenceLinker.} For each candidate triple, the system links supporting text spans and source references. This stage does not create new facts; it attaches provenance to proposed facts.

\paragraph{8. VerificationAgent.} The verification stage removes unsupported or low-quality extractions, flags hallucinations, and outputs a verified set of triples suitable for admission review.

\paragraph{9. KnowledgeOrganizer + Governance.} Each surviving triple is not inserted directly. Instead, the system executes the governance protocol:
\begin{enumerate}[nosep, leftmargin=*]
    \item route the triple to the owning domain(s),
    \item produce a \texttt{GovernanceDecision},
    \item commit the decision,
    \item append the decision to the audit log.
\end{enumerate}

This is the step that distinguishes governed creation from ordinary extraction.

\section{Deliberation, Blackboard, and Voting}
\label{sec:deliberation}

\subsection{When Deliberation Starts}

The pipeline is sequential by default. Deliberation is an exception path triggered when an agent reports low confidence. In the current implementation, a threshold around $0.7$ is used as the key decision boundary for exploratory or uncertain behavior. If an agent's confidence is above threshold, the pipeline continues. If confidence falls below threshold, the agent raises a hypothesis.

\subsection{What a Hypothesis Is}

A hypothesis is a structured claim about an extraction unit that is not trusted enough for automatic continuation. Typical hypothesis types are:
\begin{itemize}[nosep, leftmargin=*]
    \item an uncertain entity,
    \item an uncertain triple,
    \item an uncertain relation type.
\end{itemize}

Each hypothesis includes:
\begin{itemize}[nosep, leftmargin=*]
    \item the proposed content,
    \item the posting agent,
    \item the initial confidence,
    \item the supporting evidence,
    \item the current pipeline context.
\end{itemize}

\subsection{What the Blackboard Does}

The blackboard is a shared forum inside the shared memory system. It is not a chat interface; it is a structured coordination object. An agent posts a hypothesis to the blackboard, and the deliberation coordinator asks a set of evaluator agents to vote on it.

\subsection{Why Non-Domain-Specialists Still Matter}

The relevant question is not whether every voter is a topical expert on the content. The important point is that each voting agent is an \emph{expert on a criterion}. In the current extraction-time deliberation path, the three most important voters are:
\begin{itemize}[nosep, leftmargin=*]
    \item \textbf{EntityExtractor}: expert on whether the entity hypothesis is plausible under the extracted schema and text span.
    \item \textbf{RelationExtractor}: expert on whether the relation structure is coherent and bindable.
    \item \textbf{EvidenceLinker}: expert on whether the text actually supports the proposal.
\end{itemize}

This means the value of a vote comes from specialized perspective, not necessarily topical ownership. Deliberation here is about extraction reliability, not downstream domain governance. Domain governance happens later, when the triple enters the governed KG.

\subsection{Voting and Resolution}

Each voting agent evaluates the posted hypothesis and returns a vote with rationale. The coordinator aggregates the votes and resolves the hypothesis into a keep/no-keep decision. If the votes conflict materially, the coordinator can trigger a short structured debate before final resolution. The final decision and rationale are written back to shared memory, creating an audit trail for the extraction process itself.

\section{Shared Memory}
\label{sec:shared-memory}

The shared memory system has two roles in this paper.

\paragraph{First,} it supports extraction-time coordination. The blackboard, pending hypotheses, vote rationales, and local task state all live here.

\paragraph{Second,} it supports corpus-scale consistency. The memory stores canonical entities, aliases, prior extractions, and procedural patterns that help later documents reuse earlier structure rather than rebuilding it from scratch.

Conceptually, the memory system has four regions:
\begin{itemize}[nosep, leftmargin=*]
    \item \textbf{episodic memory}: per-document processing history,
    \item \textbf{working memory}: current pipeline state,
    \item \textbf{semantic memory}: accepted graph structure and alias mappings,
    \item \textbf{procedural memory}: learned extraction and organization heuristics.
\end{itemize}

\section{Output Schema}
\label{sec:schema}

The output of the creation pipeline is a JSON serialization of the governed graph. At minimum, the structure contains:
\begin{itemize}[nosep, leftmargin=*]
    \item \texttt{knowledge\_graph}
    \item \texttt{org\_chart}
    \item \texttt{governance\_mode}
    \item \texttt{audit\_log}
\end{itemize}

\subsection{Entity Schema}

Each entity stores:
\begin{itemize}[nosep, leftmargin=*]
    \item \texttt{id}: canonical identifier,
    \item \texttt{labels}: surface forms and aliases,
    \item \texttt{type}: entity type,
    \item \texttt{metadata}: provenance and confidence information.
\end{itemize}

\paragraph{Example.}
\begin{quote}
\small
\texttt{id}: ``chronic\_metabolic\_inflammation'' \\
\texttt{labels}: [``chronic metabolic inflammation''] \\
\texttt{type}: ``metabolic condition'' \\
\texttt{metadata}: \{source\_document, source\_article, confidence, source\_segment\}
\end{quote}

\subsection{Triple Schema}

Each triple stores:
\begin{itemize}[nosep, leftmargin=*]
    \item \texttt{subject}
    \item \texttt{relation}
    \item \texttt{object}
    \item \texttt{confidence}
    \item \texttt{source}
    \item \texttt{metadata}
\end{itemize}

\subsection{Org Chart}

The org chart is the organizational view of the graph. Each domain contains:
\begin{itemize}[nosep, leftmargin=*]
    \item a domain identifier,
    \item a label,
    \item a description,
    \item seed entity and relation types,
    \item the current entity membership,
    \item the domain-specific subgraph,
    \item cross-domain relations linking the domain to others.
\end{itemize}

\paragraph{Example domain.}
\begin{quote}
\small
\texttt{id}: \texttt{research\_methods} \\
\texttt{label}: ``research methods''
\end{quote}

\subsection{Org Chart Summary}

The org chart summary is the compressed textual view of the organizational structure. It lists domain names, short descriptions, topical hints, relation signatures, and entity coverage. This summary is consumed by downstream applications such as QA routing, but it is also useful analytically because it makes the domain structure legible to humans.

\section{Applications Built on Top of the Governed KG}
\label{sec:applications}

The paper's thesis is not that the system is primarily a QA architecture. Rather, the governed KG is the core object, and applications sit on top of it.

\paragraph{Question answering.} A QA orchestrator can read the org chart summary, route a question to one or more domain experts, collect subgraph-grounded answers, and synthesize a final response. This demonstrates that the governed graph is usable memory.

\paragraph{Incremental enrichment.} New documents can be added later. Candidate triples are extracted against the new documents, differenced against the base graph, and reviewed under the same governance protocol before entering memory. This demonstrates knowledge proliferation under governance.

\section{Experimental Setup}
\label{sec:setup}

We evaluate the system in two modes:
\begin{enumerate}[nosep, leftmargin=*]
    \item \textbf{Fixed-schema benchmark mode}, where the extraction pipeline is evaluated against SciERC gold annotations using hard extraction metrics.
    \item \textbf{Governance evaluation mode}, where we evaluate routing, auditability, structural coverage, and accept/reject behavior of the governed graph.
\end{enumerate}

This distinction is important. The fixed-schema benchmark measures whether the creation pipeline is competent. The governance evaluation measures the paper's main contribution.

All experiments use \texttt{gemma4:31b} at every model tier (small / medium / large). The architecture supports tiered models, but we use a single model across tiers to isolate the effect of governance from any tier-mixing effects. Extraction-time deliberation (Section~\ref{sec:deliberation}) is supported by the system but disabled for all reported runs; the runs we report use the conservative no-blackboard configuration, so the gains attributed to governance below are achieved \emph{without} the multi-agent voting path.

\section{Results}
\label{sec:results}

All numbers reported in this section are computed from artifacts in \texttt{evaluation/results/}; the headline aggregates are stored in \texttt{paper\_table\_consolidated.json} (extraction) and \texttt{qa\_corrected\_summary.json} (downstream QA). Unless otherwise noted, every run uses \texttt{gemma4:31b} at every model tier with a fixed SciERC schema, so the differences between rows reflect the governance layer rather than model size.

We organize the evaluation around three claims: (i) governance improves the graph that is admitted, (ii) governance changes \emph{which} facts are admitted rather than acting as a thinner filter on top of the same graph, and (iii) governance scales structurally --- even when the underlying extractor becomes noisier on larger corpora, the governance layer remains correct. A downstream-QA pilot then asks whether placing the governance stack between the graph and a question answerer hurts answer quality.

\subsection{Governance as an Intervention on Graph Admission}
\label{sec:exp-extraction}

To isolate the effect of governance, we extract a knowledge graph from the SciERC test split twice with the same extractor: once under \emph{flat} insertion (entities and triples are written directly into the graph) and once under \emph{governed} insertion (the same extractor paired with domain ownership, the triage governance board, and an audit trail). The two pipelines share documents, model, and schema, so any difference in their outputs is attributable to the governance layer. We repeat the comparison at 10 and 50 documents.

Extraction quality is measured against SciERC's gold annotations. \emph{Entity F1} is the standard precision-recall harmonic mean over the produced named-entity set. \emph{Triple F1} applies the same harmonic mean to $(\text{subject}, \text{relation}, \text{object})$ triples and counts a triple correct only when all three slots match the gold; it is therefore a substantially harder metric than entity F1. \emph{Triple hallucination rate} is the fraction of produced triples whose key does not appear in the gold annotations, and \emph{zero-triple docs} counts documents for which the system emitted no triples at all.

\begin{table}[t]
\centering
\small
\begin{tabular}{lccccc}
\toprule
\textbf{Scale} & \textbf{Config} & \textbf{Entity F1} & \textbf{Triple F1} & \textbf{Triple Halluc.} & \textbf{Zero-triple docs} \\
\midrule
10 docs & Flat & 0.481 & 0.079 & $87.2\%$ & 2 \\
10 docs & Governed (triage) & \textbf{0.642} & \textbf{0.292} & \textbf{$62.5\%$} & 2 \\
10 docs & $\Delta$ (governed $-$ flat) & $+33\%$ rel & \textbf{$+271\%$ rel} & $-24.7$ pp & --- \\
\midrule
50 docs & Flat & 0.633 & 0.153 & $85.8\%$ & 0 \\
50 docs & Governed (triage) & 0.629 & \textbf{0.189} & \textbf{$78.5\%$} & 3 \\
50 docs & $\Delta$ (governed $-$ flat) & $-0.4\%$ rel & \textbf{$+23.5\%$ rel} & $-7.3$ pp & --- \\
\bottomrule
\end{tabular}
\caption{Per-config extraction quality on 10 and 50 SciERC test documents (fixed schema, \texttt{gemma4:31b}). Governance preserves entity F1 while substantially lifting triple F1 and lowering triple hallucination at both scales.}
\label{tab:extraction}
\end{table}

At both scales, governance preserves entity F1 (within one percentage point of flat) while substantially lifting triple F1: a $+271\%$ relative improvement at 10 documents and $+23.5\%$ at 50 documents. Triple hallucination falls by 7--25 percentage points, and the wall-clock cost of the governance layer is small (Table~\ref{tab:governed-vs-flat-structure}). A natural concern is that this improvement reflects governance acting as a thin filter on flat's output rather than a substantive intervention; we rule this out by measuring the structural overlap between the two graphs.

The \emph{Jaccard similarity} of two sets $A$ and $B$ is $|A \cap B| / |A \cup B|$, a number between $0$ and $1$: $1.0$ indicates identical sets, $0.0$ disjoint sets, and $0.5$ roughly half overlap. We compute Jaccard separately on the entity and triple sets of the governed and flat graphs at 50 documents.

\begin{table}[t]
\centering
\small
\begin{tabular}{lcl}
\toprule
\textbf{Metric} & \textbf{Value} & \textbf{Plain reading} \\
\midrule
Entity overlap (Jaccard) & 0.677 & The two graphs largely agree on which entities to include. \\
Triple overlap (Jaccard) & \textbf{0.180} & The two graphs disagree on $\approx 80\%$ of triples. \\
Triples unique to governed & 170 & Triples governance admitted that flat never proposed. \\
Triples unique to flat & 303 & Triples flat admitted that governance rejected. \\
Cross-domain fraction & $43.4\%$ & Fraction of governed triples whose subject and object live in \\
 & & different domains; cross-domain edges are routinely admitted. \\
Wall-clock overhead & $+4.55\%$ & $+1{,}890$ s on top of the 41{,}506-s flat run. \\
\bottomrule
\end{tabular}
\caption{Structural impact of governance at 50 documents. The two graphs largely agree on entities (Jaccard $0.677$) but disagree on most edges (triple Jaccard $0.180$): governance admits 170 triples flat never produced and rejects 303 triples flat would have admitted.}
\label{tab:governed-vs-flat-structure}
\end{table}

The two graphs largely agree on entities (Jaccard $0.677$) but disagree on most edges (triple Jaccard $0.180$). Governance admits 170 triples that flat never proposed and rejects 303 triples flat would have admitted; cross-domain triples --- those whose subject and object live in different domains --- account for $43.4\%$ of the governed graph. The improvement in Table~\ref{tab:extraction} is therefore the result of a substantial change in which facts are admitted, not of post-hoc filtering. To make the governance behavior concrete, Table~\ref{tab:triage} reports the activity of the triage board on the 50-document run: each candidate triple was auto-approved (low risk), escalated to LLM review, or revised/rejected.

\begin{table}[t]
\centering
\small
\begin{tabular}{lc}
\toprule
\textbf{Route} & \textbf{Count} \\
\midrule
Auto-approved (low-risk) & 163 \\
Escalated to LLM review & 130 \\
\quad $\cdot$ Cross-domain review & 119 \\
\quad $\cdot$ Conflict review & 11 \\
Reviewer outcome: approve & 99 \\
Reviewer outcome: revise & 17 \\
Reviewer outcome: other (escalate/reject) & 14 \\
\bottomrule
\end{tabular}
\caption{Triage board activity on the 50-document run. The board exercises both auto-approval and LLM review; 31 of 130 escalations produced a revision or rejection, indicating non-trivial admission control rather than a rubber stamp.}
\label{tab:triage}
\end{table}

Of the 130 escalations, 31 produced a revision or a rejection, indicating that the board functions as a non-trivial admission gate rather than as a rubber stamp.

\subsection{Governance Structure at Scale}
\label{sec:governance-scale}

We next evaluate the governed pipeline on 100 SciERC test documents. The question here is not whether governance beats flat --- that comparison was already made --- but whether the governance layer continues to behave correctly as the underlying extractor becomes noisier with corpus size. In addition to the extraction-side metrics from Section~\ref{sec:exp-extraction}, we report a set of governance-side metrics: \emph{routing exact match} is the fraction of triples routed to the same domain a gold-standard router would choose; \emph{routing recall} and \emph{precision} are computed against expected-domain labels; \emph{cross-domain routing recall} is recall on the subset of triples that span two domains; \emph{domain coverage} is the fraction of expected domains for which at least one triple was admitted; \emph{governance completeness} is the fraction of admitted triples that traversed the governance protocol rather than a bypass path; and \emph{audit-trail integrity} is the fraction of admitted triples carrying a complete audit-trail entry.

\begin{table}[t]
\centering
\small
\begin{tabular}{lc}
\toprule
\textbf{Metric} & \textbf{Value} \\
\midrule
KG entities & 1{,}201 \\
KG triples & 338 \\
Entity F1 & 0.612 \\
Triple F1 & 0.133 \\
Orphan-entity fraction & $63.8\%$ \\
Zero-triple docs & 30 / 100 \\
\midrule
\textbf{Routing exact match} & \textbf{1.000} \\
\textbf{Routing recall (expected domain)} & \textbf{1.000} \\
\textbf{Routing precision (expected domain)} & \textbf{1.000} \\
\textbf{Cross-domain routing recall} & \textbf{1.000} \\
\textbf{Domain coverage} & \textbf{0.984} \\
\textbf{Governance completeness} & \textbf{1.000} \\
\textbf{Audit-trail integrity} & \textbf{1.000} \\
\bottomrule
\end{tabular}
\caption{Governed-only run on 100 SciERC test documents. Extraction quality drifts down at this scale (orphan fraction $63.8\%$, 30 of 100 documents emit no triples), but the governance-side metrics --- routing, completeness, audit, domain coverage --- all sit at or above $0.98$.}
\label{tab:governance-100docs}
\end{table}

Extraction quality drifts down at this scale: the orphan-entity fraction grows to $63.8\%$ and 30 of 100 documents emit no triples. The governance-side metrics, however, all sit at or above $0.98$. Routing, governance completeness, and audit-trail integrity are exact ($1.000$), and domain coverage reaches $0.984$ over the 100-document corpus. The governance layer remains correct even when the extractor underneath does not.

To probe the governance function $\gamma$ directly, we present the board with 50 probe triples --- half drawn from the gold and half corrupted versions of correct triples --- and measure whether the board admits the correct ones (\emph{accept recall}) and rejects the corrupted ones (\emph{reject recall}); the \emph{false-accept rate} counts corrupted triples that were incorrectly admitted. We compare two operating modes: \texttt{audit\_only} (the default triage policy used elsewhere in the paper) and \texttt{strict} (a conservative ceiling that prefers rejection on uncertainty).

\begin{table}[t]
\centering
\small
\begin{tabular}{lccc}
\toprule
\textbf{Mode} & \textbf{Accept recall (positive)} & \textbf{Reject recall (negative)} & \textbf{False-accept rate} \\
\midrule
Audit-only & 0.96 & 0.84 & 0.16 \\
\textbf{Strict} & 0.56 & \textbf{1.00} & \textbf{0.00} \\
\bottomrule
\end{tabular}
\caption{Strict-vs-triage review on the 100-document governed KG with 50 probe triples. Strict review attains a zero false-accept rate at the cost of rejecting $44\%$ of correct proposals; triage retains $96\%$ of correct admissions while rejecting $84\%$ of corrupted probes.}
\label{tab:strict-governance}
\end{table}

Strict review attains a zero false-accept rate at the cost of rejecting $44\%$ of correct proposals. The triage default admits $96\%$ of correct probes while still rejecting $84\%$ of corrupted ones --- the more useful trade-off when the goal is to grow the graph rather than freeze it. Together, Tables~\ref{tab:governance-100docs} and~\ref{tab:strict-governance} show that the governance layer scales structurally in the sense that matters for the contribution of this paper: routing, audit, and completeness remain exact, and the system retains a provably safe operating mode whenever it is needed.

\subsection{Downstream Question Answering on the Governed Graph}
\label{sec:qa-results}

Finally, we treat the 50-document governed graph as memory and ask whether placing the governance stack between this memory and a downstream question answerer affects answer quality. We compare four configurations on the same eight-question benchmark, covering single-hop, comparison, cross-domain, and negative questions (the latter requiring abstention when the queried fact is absent from the graph). All four configurations share the underlying KG and model and differ only in retrieval and reasoning. \texttt{flat\_path\_basic} performs flat RAG-style retrieval over the triple list. \texttt{graphrag\_basic} is a GraphRAG-style baseline that builds entity-relation communities, summarizes them, and answers from retrieved community summaries; it does not use governance routing or a critic. \texttt{domain\_basic} activates the governed domain router only, using domain ownership produced at extraction time. \texttt{domain\_advanced} adds active graph exploration, a critic, and provenance-aware memory on top of the router.

Answer quality is evaluated via atomic-fact decomposition with KG verification (KGAFE). For each answer, we extract atomic factual claims and verify each against the KG and the gold answer. \emph{Faithfulness} is the fraction of the answer's atomic claims entailed by the KG; \emph{precision / support} is the fraction with explicit support in the retrieved evidence; \emph{hallucination rate} is the fraction unsupported by the KG; and \emph{coverage} is the fraction of gold-answer atomic claims that the system's answer actually addresses. The \emph{corrected QA score} is the composite KGAFE score combining faithfulness, precision, and coverage, with one correction relative to the raw aggregate: correct abstention on a negative question is no longer scored as hallucination, since refusing to assert facts absent from the graph is the desired behavior of a governed system (Appendix~\ref{sec:qa-appendix} details the raw aggregate and the correction).

\begin{table}[t]
\centering
\small
\begin{tabular}{lccccc}
\toprule
\textbf{System} & \textbf{Corrected QA} & \textbf{Faithfulness} & \textbf{Precision} & \textbf{Halluc.\ rate} & \textbf{Coverage} \\
\midrule
\texttt{flat\_path\_basic}    & 0.691           & \textbf{1.000} & \textbf{1.000} & \textbf{$0.0\%$} & 0.938 \\
\texttt{graphrag\_basic}      & 0.673           & 0.984          & 0.984          & $1.6\%$          & 0.875 \\
\texttt{domain\_basic}        & 0.680           & \textbf{1.000} & 0.958          & $4.2\%$          & 0.938 \\
\textbf{\texttt{domain\_advanced}} & \textbf{0.698} & \textbf{1.000} & 0.991       & $0.9\%$          & \textbf{1.000} \\
\bottomrule
\end{tabular}
\caption{Four-way QA on the 50-document governed KG ($n=8$). All four configurations land in a tight band: faithfulness $\geq 0.98$ and hallucination $\leq 4.2\%$ across the board. The full governed stack (\texttt{domain\_advanced}) has the highest corrected composite score, the highest coverage ($1.000$, the only system that fully addresses every gold question), and the lowest hallucination rate among governance configurations.}
\label{tab:qa-fourway}
\end{table}

We do \emph{not} claim a statistical QA win on $n=8$. The honest reading of Table~\ref{tab:qa-fourway} is that governance does not hurt downstream QA, and is the only configuration to reach full coverage. A larger QA benchmark is left to future work.

\subsection{What the Results Mean}

Taken together, the results support five conclusions.

\begin{enumerate}[nosep, leftmargin=*]
    \item \textbf{Governance improves the admitted graph.} Triple F1 is $+271\%$ relative at 10 documents and $+23.5\%$ relative at 50 documents, with hallucination down 7--25 pp, at a $4.55\%$ wall-clock cost (Table~\ref{tab:extraction}).
    \item \textbf{Governance changes graph structure, not just volume.} Triple Jaccard of $0.180$ between governed and flat at 50 documents (Table~\ref{tab:governed-vs-flat-structure}); $303$ flat-only and $170$ governed-only triples.
    \item \textbf{Governance scales structurally.} At 100 documents, routing, completeness, audit-trail integrity, and domain coverage remain at or above $0.98$ (Table~\ref{tab:governance-100docs}) even as extraction-side noise grows.
    \item \textbf{Governance supports safe operating modes.} Strict review achieves $0\%$ false accepts; triage recovers $96\%$ of good triples while still rejecting $84\%$ of bad ones (Table~\ref{tab:strict-governance}).
    \item \textbf{Governance does not degrade downstream QA, and reaches full coverage.} On the same 50-document governed KG, all four QA configurations achieve high faithfulness ($\geq 0.98$) and low hallucination ($\leq 4\%$) on an 8-question pilot once correct abstention is properly credited (Table~\ref{tab:qa-fourway}). The full governed stack (\texttt{domain\_advanced}) is the unique system at full coverage ($1.000$) with the lowest hallucination rate among the governance configurations. We do not claim a statistical QA win on this small benchmark; the QA result is supporting evidence that the governance layer extends coverage without harming faithfulness.
\end{enumerate}

\section{Limitations}
\label{sec:limitations}

We report results on a single extraction benchmark (SciERC), with a single model family (\texttt{gemma4:31b}) used at every tier, at scales up to 100 documents. Extending to additional benchmarks (e.g., DocRED, CoNLL) and to mixed-tier model configurations is left to future work. The downstream QA evaluation is a small pilot ($n=8$ for the four-way comparison; $n=5$ for the judge-mode rerun) on the 50-document governed KG, and the QA-quality differences across configurations are within noise on this benchmark; we therefore report QA as supporting evidence that governance does not degrade downstream answering, not as a statistical win. The extraction-time deliberation path described in Section~\ref{sec:deliberation} is supported by the system but disabled in all runs reported here; ablating deliberation on/off is left to future work, and the governance gains we report do not depend on it.

\section{Related Work}
\label{sec:related}

\paragraph{Multi-agent KG construction.} Recent work such as KARMA~\cite{karma2025} studies multi-agent role specialization for knowledge graph construction. Our work differs in emphasis: the central object is not the extraction pipeline alone, but a governed graph with explicit ownership, admission control, and auditability. We do not position the system as a state-of-the-art extractor and we do not run head-to-head comparisons against fine-tuned IE systems; the contribution we evaluate is the governance layer over the extracted graph.

\paragraph{Information extraction benchmarks.} SciERC~\cite{scierc2018} remains a standard benchmark for scientific entity and relation extraction. We use SciERC in fixed-schema mode to establish extraction competence, but do not position our contribution as a new state-of-the-art extractor.

\paragraph{Factual evaluation.} FActScore~\cite{factscore2023} and SAFE~\cite{safe2024} evaluate factuality of generated text. Our work is adjacent in spirit but structurally different: we evaluate the creation and governance of machine knowledge itself rather than only the factuality of generated answers.

\paragraph{LLM-based scientific systems.} SciAgents~\cite{sciagents2024} and related systems show that role-specialized LLM agents can structure complex workflows. We adopt this general design pattern but tie it to a formal governed memory object and a creation-time governance protocol.

\section{Conclusion}

We introduced the Governed Knowledge Graph, a knowledge representation in which extracted facts are organized under domain-owned subgraphs and admitted through an explicit governance function. The system couples a multi-agent creation pipeline with domain bootstrap, ownership assignment, proposal/commit semantics, and an audit trail over admitted triples. Empirically, the strongest evidence is governance-specific rather than extraction-specific: exact routing on the created graph, near-complete domain coverage, complete auditability, and strong strict-review behavior. These results support the view that governed memory is a valid and useful object of study in its own right, with question answering and enrichment acting as applications built on top of it.

\bibliographystyle{plainnat}

\begin{thebibliography}{10}

\bibitem[Karma et~al.(2025)]{karma2025}
KARMA: Knowledge Graph-Based Agent for Reasoning and Multi-Agent Collaboration.
\newblock In \emph{NeurIPS}, 2025.

\bibitem[Luan et~al.(2018)]{scierc2018}
Y.~Luan, L.~He, M.~Ostendorf, and H.~Hajishirzi.
\newblock Multi-Task Identification of Entities, Relations, and Coreference for
  Scientific Knowledge Graph Construction.
\newblock In \emph{EMNLP}, 2018.

\bibitem[Min et~al.(2023)]{factscore2023}
S.~Min, K.~Krishna, X.~Lyu, M.~Lewis, W.~Yih, P.~Koh, M.~Iyyer, L.~Zettlemoyer, and H.~Hajishirzi.
\newblock FActScore: Fine-grained Atomic Evaluation of Factual Precision in Long Form Text Generation.
\newblock In \emph{EMNLP}, 2023.

\bibitem[Wei et~al.(2024)]{safe2024}
J.~Wei, C.~Anil, et~al.
\newblock Long-form factuality in large language models.
\newblock \emph{arXiv preprint arXiv:2403.18802}, 2024.

\bibitem[Ghafarollahi and Buehler(2024)]{sciagents2024}
A.~Ghafarollahi and M.~J.~Buehler.
\newblock SciAgents: Automating scientific discovery through multi-agent intelligent graph reasoning.
\newblock \emph{arXiv preprint arXiv:2409.05556}, 2024.

\end{thebibliography}

\appendix

\section{QA Scoring Correction}
\label{sec:qa-appendix}

The QA aggregate reported in Table~\ref{tab:qa-fourway} is the \emph{corrected} aggregate. The raw KGAFE composite penalized correct abstention on negative questions: when the gold answer was ``the knowledge graph does not contain a relationship between $X$ and $Y$,'' a model answer of the form ``there is no evidence that $X$ relates to $Y$'' was decomposed into a single atomic fact that could not be verified against any KG triple, and was scored as \texttt{hallucination}. Because all four configurations correctly abstain on negative questions, this artifact pushed every system's raw hallucination rate up by a fixed amount and inflated the apparent KGAFE-composite spread. The corrected aggregate (stored in \texttt{evaluation/results/qa\_corrected\_summary.json}) credits correct abstention on negative questions and is the appropriate signal for governance: abstention on a triple that does not exist is the desired behavior of the governance layer.

\end{document}
